\section{Introduction}
Procedural memory gives an agent a way to reuse a successful method without deriving it again. Yet success is conditional. A length-delimited record may have been complete before it reached a parser; a cache may have belonged to one request; a diagnostic object may have contained only public fields. When the target changes these assumptions, the remembered procedure can remain an accurate account of the source task while becoming an incomplete guide to the new one.

We ask: \emph{when a coding agent receives a source-valid procedure in a target context where its assumptions no longer hold, how does that memory change functional completion and focal security failures?} A second question asks whether a generic reminder to check applicability changes that behavior. The intervention is deliberately small: supply a frozen procedural memory, supply a matched irrelevant memory, append an applicability reminder, or supply no memory. The target task and its changed assumption are held fixed across these conditions.

Our original hypothesis was that source-correct memory would increase functional completions that violate the target's focal security property. The completed experiment does not support this as a general effect. Across six configurations, the corrected primary contrast has both positive and negative point estimates, with substantial uncertainty. This result does not make the intervention behaviorally inert. Matched cases reveal unchecked reuse in some runs, recognition of source limitations in others, and failures to complete the task that would be mistaken for security improvements if functionality were omitted.

The study makes three contributions. First, it operationalizes source-valid but target-inapplicable procedural memory in thirteen controlled synthetic families with separate functionality and security witnesses. Second, it reports corrected, within-configuration treatment comparisons and the full joint outcome distribution, preserving technical invalidity as missing measurement. Third, it connects selected visible transcripts to saved implementation changes and evaluated outcomes, showing why neither a witness failure nor silence about memory establishes how a procedure was used. These are behavioral case studies selected after outcome inspection, not estimates of behavioral prevalence.

The scope is intentionally specific. All experimental targets already contain a context shift; there is no randomized unshifted target arm. The design estimates the effect of supplying memory under those conditions, not the effect of shift itself. Memories are supplied directly, so retrieval quality, memory selection, and online memory formation are outside the intervention. The synthetic setting supports controlled comparisons but does not establish performance in real repositories.
